
\clearpage
\chapter{מערכות ניתוב כלים מוקדמות}

\section{\textenglish{ReAct}: הצעד הראשון}

\textenglish{ReAct} \cite{yao2022react} הגדיר פרוטוקול פשוט: חשיבה (\textenglish{Thought}) $\to$ פעולה (\textenglish{Action}) $\to$ תצפית (\textenglish{Observation}) $\to$ חשיבה מחדש. כל שלב מיוצג כטקסט ב-\textenglish{prompt} הנשלח למודל. ניסויים על \textenglish{HotpotQA} ו-\textenglish{FEVER} הדגימו שיפור של \textenglish{34\%} על \textenglish{CoT} בלבד.

\section{\textenglish{Toolformer}: למידת שימוש בכלים}

\textenglish{Toolformer} \cite{schick2023toolformer} הציג גישה שונה לחלוטין: במקום לתכנת ידנית מתי להשתמש בכלי, המודל לומד לבד להוסיף קריאות \textenglish{API} בתוך הטקסט שהוא מייצר. האימון משתמש ב\textenglish{-self-supervised} הסינון: קריאות \textenglish{API} שמפחיתות \textenglish{perplexity} על הטקסט הבא נשמרות; אחרות נמחקות.

\section{\textenglish{ToolLLM}: מאגר כלים בקנה מידה גדול}

\textenglish{ToolLLM} \cite{qin2023toolllm} הרחיב את הגישה ל-\textenglish{16,000+} כלי \textenglish{API} אמיתיים מ-\textenglish{RapidAPI}. ה-\textenglish{DFSDT} (\textenglish{Depth-First Search-based Decision Tree}) מאפשר לסוכן לחקור אפשרויות כלים בצורה שיטתית ולהתאושש מכשלים:

\begin{equation}
\text{DFSDT}(s_t, \mathcal{T}) = \arg\max_{a \in \mathcal{A}} Q(s_t, a; \mathcal{T})
\end{equation}

כאשר \(\mathcal{T}\) היא עץ ההחלטות ו-\(Q\) הוא ציון האיכות המשוער מהמודל.
